% =====================================================================
%  Geneigtes Sägeblatt (Vorderansicht, Blick der Bedienperson): Das Blatt
%  neigt sich nach rechts. Den Längsanschlag dann links vom Blatt setzen,
%  damit kein Reststück unter dem Blatt eingeklemmt wird.
%  Selbst gezeichnet (TikZ), Lizenz: CC BY-SA 3.0
% =====================================================================
\begin{tikzpicture}[x=1mm, y=1mm, font=\scriptsize]

  % Tisch mit um 30° geneigtem Blatt, Drehpunkt auf der Tischoberfläche
  \newcommand{\neigetisch}{%
    \fill[gehaeusegruen!70] (6,4) rectangle (76,12);
    \draw[tisch] (6,12) rectangle (76,14);
    \draw[hilfslinie, dashed] (41,14) -- (41,36);}
  \newcommand{\neigeblatt}{%
    \draw[blatt, rotate around={-30:(41,14)}] (40.4,14) rectangle (41.6,34);}

  % ---------------- 1 richtig --------------------------------------
  \feld{0}{0}{82}{52}{1\quad Anschlag links vom Blatt}
  \neigetisch
  \draw[werkstueck] (18,14) -- (41,14) -- (45.6,22) -- (18,22) -- cycle;
  \draw[werkstueck] (41.9,14) -- (60,14) -- (60,22) -- (46.5,22) -- cycle;
  \neigeblatt
  \draw[metall] (12,14) rectangle (18,26);
  \draw[drehung] (41,30) arc[start angle=90, end angle=60, radius=16];
  \node at (50,37) {0--45°};
  \node[anchor=west, align=left] at (60,30) {Blatt neigt\\sich vom\\Anschlag weg};
  \node[anchor=south] at (15,26) {Anschlag};
  \pic at (74,45) {ok};

  % ---------------- 2 falsch ---------------------------------------
  \begin{scope}[shift={(88,0)}]
    \feld{0}{0}{82}{52}{2\quad Anschlag rechts vom Blatt}
    \neigetisch
    \draw[werkstueck] (14,14) -- (41,14) -- (45.6,22) -- (14,22) -- cycle;
    % Reststück unter dem geneigten Blatt eingeklemmt
    \draw[werkstueck] (41,14) -- (51,14) -- (51,22) -- (45.6,22) -- cycle;
    \draw[metall] (51,14) rectangle (57,26);
    \neigeblatt
    \draw[kraft, nokrot] (48,24) -- (48,34);
    \node[nokrot, anchor=west, align=left] at (58,32) {Reststück\\klemmt unter\\dem Blatt};
    \pic at (74,45) {nok};
  \end{scope}
\end{tikzpicture}
